\section{Results}\label{sec:results}

%Presents the results of the procedure and includes:
%1) all data collected in tabular form if possible;
%2) all required graphs;
%3) all results in a logical sequence;
%4) objective observations;
%Note: all figures, tables and graphs required should include
%titles, labels, units and legends where appropriate.

%Present processed data, not raw data.

\begin{figure}[htb]
    \centering
    \includegraphics[width=\linewidth]{figures/magnet_strength}
    \caption{Magnetic field strength of the long and short magnets as a function of inverse cube distance. }
    \label{fig:method1}
\end{figure}

In Figure ~\ref{fig:method1}, the magnetic field strength was plotted against inverse cube distance in order to
approximate the value of the dipole moment.
As the dipole moment equation is only a far field approximation, the theoretical linear relationship corresponds to
the data collected the furthest away from the magnet.

\begin{figure}[htb]
    \centering
    \includegraphics[width=\linewidth]{figures/short_fall}
    \caption{EMF induced by releasing the magnetic from rest.}\label{fig:short_fall}
\end{figure}

\begin{figure}[htb]
    \centering
    \includegraphics[width=\linewidth]{figures/long_fall}
    \caption{EMF induced by releasing the long magnet from rest.}\label{fig:long_fall}
\end{figure}

In Figures ~\ref{fig:short_fall} and ~\ref{fig:long_fall}, the EMF is induced from the flux of the magnetic field
lines changing with time as the magnet falls.
The voltage peak-peak amplitude increases with each successive coil the magnet falls through.

By performing a Riemann sum over the positive and negative values, we can resolve the flux through each coil.
Table ~\ref{tab:riemann1} and ~\ref{tab:riemann2} shows the calculated result from the falling short magnet.
Each coil has close agreement between their positive and negative phases, resulting in approximately equal and
opposite total flux.

\begin{table}[htb]
    \caption{Calculated Flux Values - Positive Phases}
    \centering
\begin{tabular}{lr}
    \toprule
Flux Component & Quantity {[}Wb{]} \\
    \midrule
fluxA          & -1.117e-04        \\
fluxB          & -1.081e-04        \\
fluxC          & -1.133e-04        \\
fluxD          & -1.104e-04        \\
flux total     & -2.218e-05        \\
    \bottomrule
\end{tabular}\label{tab:riemann1}
\end{table}

\begin{table}[htb]
    \caption{Calculated Flux Values - Negative Phases}
    \centering
\begin{tabular}{lr}
    \toprule
Flux Component & Quantity [Wb] \\
    \midrule
fluxA          & 1.106e-04        \\
fluxB          & 1.073e-04        \\
fluxC          & 1.126e-04        \\
fluxD          & 1.095e-04        \\
flux total     & 2.200e-05        \\
    \bottomrule
\end{tabular}\label{tab:riemann2}
\end{table}



\begin{figure}[htb]
    \centering
    \includegraphics[width=\linewidth]{figures/position_short}
    \caption{Z position as a function of time.}\label{fig:position_short}
\end{figure}

\begin{figure}[htb]
    \centering
    \includegraphics[width=\linewidth]{figures/position_long}
    \caption{Z position as a function of time.}\label{fig:position_long}
\end{figure}

In Figures ~\ref{fig:position_short} and ~\ref{fig:position_long}, a quadratic fit is applied to the position data
at the points where the magnet passes through the coils.
By extracting the value from the quadratic fit line, we determine that the acceleration due to gravity was
$\qty{9.41}{\meter\per\second^2}$ and $\qty{9.48}{\meter\per\second^2}$ for the short and long magnets,
respectively.
To derive an expression for speed, $v(t)$, we take the derivative of $z(t)$:

\begin{equation}
    v(t) =
    \begin{cases*}
        9.41t + 3.54 & short magnet \\
        9.48t + 3.46 & long magnet
    \end{cases*}
    \label{eq:velocity_cases}
\end{equation}

\begin{figure}[tb]
    \centering
    \includegraphics[width=\linewidth]{figures/short_coilB}
    \caption{Experimental and calculated EMF curve for short magnet.}\label{fig:short_coilB}
\end{figure}


\begin{figure}[tb]
    \centering
    \includegraphics[width=\linewidth]{figures/long_coilB}
    \caption{Experimental and calculated EMF curve for long magnet.}\label{fig:long_coilB}
\end{figure}

In Figures ~\ref{fig:short_coilB} and ~\ref{fig:long_coilB}, we compare the experimental and calculated EMF for
coil.

The calculated EMF response from the long magnet deviates more significantly compared with the short magnet.


%\clearpage


%\textcolor{red}{
%    Uncertainties in experimental values were calculated using the aid of the \textit{uncertainties} python package.
%}
